\scanpage{23}

\uline{En résumé} Dans la \uline{catégorie abélienne} \inserted{$k$-lin. $C$} de tous les $k$-modules $\frac12$ simpliciaux, ceux qui sont \inserted{str. $\partial$-}parfaits forment une sous-catégorie \uline{additive} $C_0$ \deleted{\unsure{engendr.}}, dans laquelle on a un complexe de cochaînes

\diagram{\begin{tikzcd}[ampersand replacement=\&,column sep=small]
0\arrow[r,dash] \& \Phi_{\ast}^{(0)}\arrow[r,dash] \& \Phi_{\ast}^{(1)}\arrow[r,dash] \& \Phi_{\ast}^{(2)}\arrow[r] \& \cdots\arrow[r] \& \Phi_{\ast}^{(i)}\arrow[r,"\lambda_i"] \& \Phi_{\ast}^{(i+1)}
\end{tikzcd}}

satisfaisant à ceci \par\textcircled{a} Complexe tot. acyclique dans $C$, \deleted{$\unsure{\lambda_i\ne0\ (i\geq0)}$}

\textcircled{b} Pour tout $i,j\geq0$, on a

\[\operatorname{Hom}(\Phi_{\ast}^{(i)},\Phi_{\ast}^{(j)})\simeq
\begin{cases}0\ \text{si }j\unsure{=|\ne}i,\,i+1\\
\simeq k\ (\text{engendré par id})\ \text{si }j=i\\
\simeq k\ (\text{----- }\lambda_i)\ \text{si }j=i+1\end{cases}\]

\deleted{\textcircled{c}} \margin{conséquences de \textcircled{a} \textcircled{b}}% slanted
Posant $\Psi_{\ast}^{(i)}=Z^{i+1}(\Phi^{\unsure{\ast|\bullet}})=\operatorname{Ker}_{C}(\Phi_{\ast}^{i+1}\xrightarrow{\lambda_{i+1}}\Phi_{\ast}^{i+2})$\inserted{$\simeq\operatorname{Coker}_{C}(\Phi^{i-1}\xrightarrow{\lambda_{i-1}}\Phi^i)$} \deleted{on a donc Hom $(\Psi^i$}

\[ \Psi^0\simeq\Phi^0 \]

\diagram{\begin{tikzcd}[ampersand replacement=\&,column sep=small,row sep=small]
(*)\quad 0\arrow[r,dash] \& \Psi^{i-1}\arrow[r,hook,"\unsure{\iota_i|v_i|\partial_i}"] \& \Phi^i\arrow[r,two heads,"u_i"]\arrow[dr,"\lambda_i"'] \& \Psi^i\arrow[r]\arrow[d,hook] \& 0 \quad\text{ext. can.}\\
\& \Phi^{i-1}\arrow[u,two heads,"u"']\arrow[ur,"\lambda_{i-1}"] \& \& \Phi^{i+1} \&
\end{tikzcd}}

\[\operatorname{Hom}(\Psi^i,\Psi^j)\ (\subset\operatorname{Hom}(\Phi^i,\Phi^{j+1}))=
\begin{cases}0&\text{si }j\ne i\deleted{\unsure{,i+1}}\\k&\text{si }j=i\ (\text{base id})\end{cases}\]

\[\operatorname{Hom}(\Phi^i,\Psi^j)\deleted{\unsure{\simeq}}[\subset\operatorname{Hom}(\Phi^i,\Phi^{j+1})]\simeq
\begin{cases}0&\text{si }j\ne i\\k&\text{si }j=\unsure{0|i|\theta}\ (\text{base }u_i)\end{cases}\]

\[\operatorname{Hom}(\Psi^{i\deleted{\unsure{+1}}},\Phi^j)\ (\subset\operatorname{Hom}(\Phi^i,\Phi^j))\simeq
\begin{cases}0&\text{si }j\ne i+1,\ (i\ne0)\\k&\text{si }j=i+1\ (\text{base }\unsure{v_i|\iota_i|\partial_i})\\ k&\text{si }j=i=0\ (\text{base id})\end{cases}\]

% slanted
\margin{\unsure{zéro car forcé|engendré} par le comp. \unsure{du hom} \unsure{canonique} $\Phi(i,n)\longrightarrow$ $\Psi^i\longrightarrow\unsure{\Phi^{i+1}|\Phi(i+1,m)}$}

d) Soit, \deleted{\unsure{$\Phi(i,n)$}} pour $n\in k-0$\inserted{et $i\geq1$}, $\Phi(i,n)$ l'ext. de $\Psi^i\deleted{\unsure{i}}$ par $\Psi^{i-1}$ déduite de l'ext. $\Phi^i$ $(*)$ par $n\,\mathrm{id}_{\Psi^{i-1}}$, de sorte qu'on a un hom de suites exactes

\diagram{\begin{tikzcd}[ampersand replacement=\&,column sep=small]
0\arrow[r] \& \Psi^{i-1}\arrow[r,"v_i"]\arrow[d,"n\,\mathrm{id}"'] \& \Phi^i\arrow[r,"u_i"]\arrow[d] \& \Psi^i\arrow[r,dash]\arrow[d,"\mathrm{id}"] \& 0\\
0\arrow[r] \& \Psi^{i-1}\arrow[r] \& \Phi(i,n)\arrow[r,dash] \& \Psi^i\arrow[r,dash] \& 0
\end{tikzcd}}

(\uline{NB}\quad $\Phi(i,1)=\Phi^i$) \inserted{On a aussi}

\diagram{\begin{tikzcd}[ampersand replacement=\&,column sep=small,row sep=small]
0\arrow[r,dash] \& \Psi^{i-1}\arrow[r,dash]\arrow[d,Rightarrow,"\mathrm{id}"'] \& \Phi(i,n)\arrow[r,dash]\arrow[d] \& \Psi^i\arrow[r]\arrow[d,"n\,\mathrm{id}"] \& 0\\
0\arrow[r,dash] \& \Psi^{i-1}\arrow[r] \& \Phi^i\arrow[r] \& \Psi^i\arrow[r] \& 0
\end{tikzcd}}

Alors on a

\[\operatorname{Hom}(\Phi(i,n),\Phi(j,m))\longrightarrow\operatorname{Hom}(\Psi^{i-1},\unsure{\Phi(j)|\Psi^j})\]

et si $j\ne i\unsure{+1|-1}$, ce dernier Hom est nul, d'où si $j\ne i\unsure{+1|-1}$, la structure d'extension est respectée par tt hom $\alpha:\Phi(i,n)\longrightarrow\Phi(j,m)$

\diagram{\begin{tikzcd}[ampersand replacement=\&,column sep=small,row sep=small]
0\arrow[r,dash] \& \Psi^{i-1}\arrow[r]\arrow[d] \& \Phi(i,n)\arrow[r]\arrow[d] \& \Psi^i\arrow[r,dash]\arrow[d] \& 0\\
0\arrow[r] \& \Psi^{j-1}\arrow[r,dash] \& \Phi(j,m)\arrow[r,dash] \& \Psi^j\arrow[r,dash] \& 0
\end{tikzcd}}

si \inserted{de plus} $j\ne i$, les flèches entre les extrêmes sont nulles, donc on trouve

\[\operatorname{Hom}(\Phi(i,n),\Phi(j,m))\xleftarrow{\sim}\operatorname{Hom}(\Psi^i,\Psi^{j-1})\quad\begin{matrix}\Vert\\0\ \text{si }j\ne i+1\end{matrix}\]

\uline{Cas $j=i+1$} \quad \textquotedblleft\quad \unsure{\textquotedblright} $\simeq k$

\uline{Cas $j=i$}. Tout hom induit une $p\,\mathrm{id}$ sur $\Psi^{i-1}$ et $q\,\mathrm{id}$ sur $\Psi^{i}$,

\diagram{\begin{tikzcd}[ampersand replacement=\&,column sep=small,row sep=small]
0\arrow[r,dash] \& \Psi^{i-1}\arrow[r,dash]\arrow[d,"p\,\mathrm{id}"'] \& \Phi(i,n)\arrow[r,dash]\arrow[d,"\alpha"] \& \Psi^i\arrow[r]\arrow[d,"q\,\mathrm{id}"] \& 0\\
0\arrow[r,dash] \& \Psi^{i-1}\arrow[r] \& \Phi(i,m)\arrow[r,dash] \& \Psi^i\arrow[r,dash] \& 0
\end{tikzcd}}

$p,q$ unique\textsuperscript{mt} déterminés par $\alpha$. On a alors (pour qu'il existe $\alpha$ définissant $p$ et $q$) que l'élément de $\operatorname{Ext}^1(\Psi^i,\Psi^{i-1})$ défini par $q^*(e_{j,m})-p_*(e_{i,n})$ (où $e_{i,n}$, $e_{j,m}$ définissent resp. les extensions $\Phi(i,n)$, $\Phi(j,m)$)

% archive sheet number: 22
